\documentclass[12pt]{article}

\usepackage{tgtermes}
\usepackage{epsf}
\usepackage{epstopdf}
\usepackage{amsmath}
\usepackage{graphicx}
\usepackage{booktabs}
\usepackage[colorlinks=true,linkcolor=blue,citecolor=blue]{hyperref}
\usepackage{dcolumn}
\usepackage{amsmath, amsthm, amssymb}
\usepackage{mwe}
\usepackage{url}
%\usepackage{harvard}
\usepackage{fancyheadings}
\usepackage{longtable}
\usepackage{authblk}
\usepackage{setspace}
%\usepackage[nomarkers]{endfloat}
\usepackage{float}
\usepackage{bbm}
%\usepackage{titling}
\usepackage{subcaption}
\usepackage{algorithm}
\usepackage{algorithmic}
\usepackage{import}
\usepackage[backend=biber,style=authoryear,
sorting=nyt,citestyle=authoryear]{biblatex}
\addbibresource{references.bib}
%\usepackage[nomarkers,nofiglist,notablist]{endfloat}
\usepackage{subcaption}
\usepackage{caption}

\onehalfspacing
\textwidth 6.5in \oddsidemargin 0in \evensidemargin -0.6in
\textheight 8.5in \topmargin -0.2in

\newcolumntype{L}[1]{>{\raggedright\let\newline\\
		\arraybackslash\hspace{0pt}}m{#1}}
\newcolumntype{C}[1]{>{\centering\let\newline\\
		\arraybackslash\hspace{0pt}}m{#1}}
\newcolumntype{R}[1]{>{\raggedleft\let\newline\\
		\arraybackslash\hspace{0pt}}m{#1}}
\newcolumntype{P}[1]{>{\raggedright\tabularxbackslash}p{#1}}

\newtheorem{theorem}{Theorem}[section]
\newtheorem{corollary}[theorem]{Corollary}
\newtheorem{proposition}[theorem]{Proposition}
\newtheorem{lemma}[theorem]{Lemma}

\captionsetup{justification=centering,singlelinecheck=false}


\newcommand{\xsub}[1]{%
	\mbox{\scriptsize\begin{tabular}{@{}c@{}}#1\end{tabular}}%
}

%\renewcommand{\thetable}{\Roman{table}}

\begin{document}
	
	
	
	
	\linespread{1.2}\title{\vspace{-0.5in} Exploring Nonprofit Hospital Executive Teams in the US} 
	
	\date{\today}
	
	\author{\vspace{10mm}Hanna Glenn\footnote{School of Economics, University of Queensland, hanna.glenn@uq.edu.au} }
	
	\maketitle
	%\setlength{\droptitle}{-10pt}
	
	\vspace{-0.2in}
	
	\singlespacing\maketitle


 \vspace{3mm}
	
\begin{abstract}
\small{}
\end{abstract}
	
	
	
	
	

	
	\onehalfspacing
	
	\newpage

    \section{Introduction}

    Paragraph 1: Leaders are important! Give statistics from general media that show this. 

    Executives are typically considered important drivers of a firm's practices and performance, as they are heavily involved in the culture, visions, and strategies of the firm. Public discourse emphasizes this point, stating that the highest level of leadership drives both firm and general policy landscapes (\cite{cecchi2024leadership}; \cite{boyles2023organizational}). Additionally, academic researchers have documented correlation between executive characteristics and firm outcomes, starting with \citeauthor{bertrand2003managing} (\citeyear{bertrand2003managing}), who show that manager fixed effects are an important driver of many firm behaviors. Other studies have built on this finding, showing that female executives are correlated with female employee wages and corporate strategies (\cite{flabbi2019female}; \cite{matsa2013female}), and that CEO characteristics are correlated with merger activity, culture, and performance (\cite{levi2010deal}; \cite{benmelech2015military}; \cite{graham2022corporate}; \cite{kaplan2012ceo}; \cite{custodio2013generalists}; \cite{adams2018director}; \cite{frydman2019rising}). Despite limited causal evidence, the correlation between executive leadership and firm outcomes is well established.


    \medskip

    Paragraph 2: Despite the importance of hospitals and nonprofit, we don't know as much about leadership in this context. 

    Most of this literature focuses on large, publicly traded firms, reflecting both their economic importance and the availability of detailed data in these settings.\footnote{Examples of typical data sources are BoardEx, ExecuComp, and SEC EDGAR filings.} As a result, we know less about executives in industries dominated by smaller firms or alternative ownership structures. Hospitals are one such industry with significant economic and public welfare implications. Hospital spending consistently makes up almost a third of all health care spending in the US (\cite{nlm_health_statistics}), and approximately 7 percent of people have one or more hospital stays each year (\cite{cdc2020hospitalstay}). Despite the importance of hospital outcomes and behavior, we have a much more limited understanding of who is at the helm. \citeauthor{munoz2025managers} (\cite*{munoz2025managers}) exploit a policy in Chile that increases competition in CEO recruitment and show that higher-quality CEOs reduce mortality. In contrast, \citeauthor{janke2019impact} (\cite*{janke2019impact}) show that top leaders have little impact on hospital performance in England. The only study to my knowledge directly investigating leaders in health care in the US focuses on physician practices as opposed to hospitals (\cite{la2023impact}). Other work emphasizes management practices instead of managers (\cite{bloom2015impact}; \cite{bloom2020healthy}), leaving evidence surrounding hospital executives sparse and mixed depending on the setting being considered. 
    
    


    \medskip

    Paragraph 3: Talk about what I do in the paper. 

    This paper provides the first systematic characterization of the executive landscape in hospitals in the US. I construct a novel panel of nonprofit hospital executives from 2009 to 2023 and document key features of leadership teams, including gender composition, team size, clinical experience, compensation, and turnover. I focus on nonprofit hospitals, which make up 50\% of all hospitals in the US, and staff, on average, 207 beds (\cite{ASPE_2023}). I merge these data with the American Hospital Association (AHA) Annual Survey to examine how leadership characteristics correlate with hospital outcomes. Finally, I estimate a descriptive event study to explore how changes in executive composition relate to hospital outcomes and performance.
    
 

    \medskip

    Paragraph 4: Explain any trends that I find

    \medskip

    Paragraph 5: Talk about policy implications/future research

    
    \section{Data}

    The novelty of this study lies primarily in the extraction of US-based individual executives of nonprofit hospitals from publicly available Tax Forms 990s. A Tax Form 990 is an annual informational return that tax-exempt organizations file with the Internal Revenue Service (IRS) to report their financial activities, governance, and compensation, providing transparency to regulators and the public. The tax forms have been used in countless studies across disciplines to assess nonprofit finances, program activities, and broader patterns of behavior  (\cite{froelich2000financial}; \cite{ely2023efiling}) The only other study that I know of to use it for key hospital-level personnel is \citeauthor{brickley2010board} (\cite*{brickley2010board}), who focus on a sample of approximately 300 hospitals. 
    
    I present a blank version of the form section that I focus on in Figure \ref{fig:blank_form}.\footnote{Additionally, the entire form can be viewed at \href{https://www.irs.gov/pub/irs-pdf/f990.pdf}{https://www.irs.gov/pub/irs-pdf/f990.pdf}.} Over the sample period, this particular form has changed very little apart from the addition of average hours worked per week at related organizations. In 2009, only average hours at the specified organization were required. The resulting data includes all items in table: name and title, average hours per week at specified organization and related organization, position, reportable compensation, reportable compensation from related organizations, and estimated amount of other compensation. 

    \begin{figure}[ht!]
        \centering
        \includegraphics[width=0.75\linewidth]{Paper/Objects/blank_form.png}
        \caption{Blank Form 990 Example, 2025}
        \label{fig:blank_form}
    \end{figure}

    \subsection{Identifying Hospitals and Linking Across Data Sources}\label{sec:identify}

    The first step in collecting comprehensive hospital data is to identify a list of firms (Employer Identification Numbers byt the IRS) that (1) operate hospital facilities and (2) can be linked to other hospital information, which mitigates the future steps of the data collection process. First, I download the universe of Tax Form 990 XMLs from the Internal Revenue Service (IRS).\footnote{https://www.irs.gov/charities-non-profits/form-990-series-downloads} Though the IRS only publishes the files from 2016 onward, this allows me to algorithmically identify firms that complete Schedule H of the Tax Form 990, s section that is only completed by firms that operate hospital facilities.\footnote{The implication of this process is that I do not collect data on hospitals that closed prior to 2016.} I save this list of 2,559 Employer Identification Numbers (EIN) along with their location and organization name for the next step in data collection.

    I link the IRS-based data set of EINs to hospital identifiers from the American Hospital Association (AHA) using a multi-step cleaning and fuzzy matching procedure. I begin by standardizing organization names and address fields in both datasets. For the IRS data, I construct a unified business name variable by combining multiple name fields and applying a series of cleaning rules to handle common reporting issues, including ``doing business as'' (DBA) designations, group returns, and formatting inconsistencies. I convert all text fields to uppercase, remove punctuation and common suffixes (e.g., ``INC''), and standardize addresses and ZIP codes. I apply analogous cleaning steps to the AHA data, including standardizing abbreviations (e.g., converting ``Med'' to ``Medical'') and normalizing address and location fields.

    I then perform a fuzzy string match between the IRS and AHA datasets using the Jaro-Winkler distance metric applied to organization names, restricting matches to those within a specified similarity threshold. To improve match quality, I impose additional constraints based on geographic consistency, requiring matches to occur within the same state and rewarding agreement in city, ZIP code, and address similarity. I construct a match quality index that combines these criteria, including penalties for large string distances and bonuses for exact or near-exact matches. For each EIN-year observation, I retain the match with the highest index and discard low-quality matches below a minimum threshold.

    I refine the matched sample by resolving duplicate matches and enforcing a one-to-one mapping between EINs and AHA hospital identifiers. I drop observations with multiple conflicting matches and manually verify a random sample of matches, finding a high rate of accuracy. I then augment the matched set by conducting a secondary fuzzy match based on address fields for previously unmatched observations. Finally, I manually assign identifiers for a small number of remaining unmatched EINs where reliable matches can be identified.

    This procedure yields a final crosswalk linking 1,836 EINs to AHA hospital identifiers, which I use to merge IRS filing data with hospital-level characteristics. The resulting matched sample provides a consistent mapping between nonprofit tax filings and hospital entities across datasets.


    \subsection{XML Data Collection and Parsing}

    I programmatically download Form 990 XML filings for each hospital EIN in the final list created in Section \ref{sec:identify} using ProPublica’s API and organization pages. The collection procedure proceeds in three steps. First, I normalize EINs to a standardized nine-digit numeric format and query the ProPublica search API for each EIN. I retain only cases where the search returns exactly one matching organization, excluding EINs that produce zero or multiple matches to avoid ambiguity. Second, for each validated EIN, I scrape the corresponding ProPublica organization page to identify all available XML filings. I collect hyperlinks corresponding to XML downloads by selecting links labeled ``XML'' or links whose URLs contain patterns consistent with XML files (e.g., ``download-xml'' or URLs ending in ``.xml''). Third, I download all available XML files for each EIN. 

    I parse each XML file to construct both firm-level and person-level datasets. Because Form 990 XML schemas vary across filing types and over time, I implement a schema-agnostic parser that searches across multiple possible node structures. For each XML file, I also store diagnostic information on the parsing process, including the schema used for extracting person-level data, the number of nodes identified, the number of individuals extracted, and indicators for missing names, titles, and firm-level fields. I flag observations for review if firm-level information is incomplete, if no individuals are extracted, or if extraction relies on text-based fallback methods. I also log parsing errors to ensure transparency and reproducibility.

    \subsubsection{Firm-Level Data}

    From each XML file, I extract the following information: EIN and organization name, location (city, state, ZIP code, and country), filing year, and return type (e.g., Form 990, 990-EZ, 990-PF).\footnote{Form 990 filings vary by organization size and type: the standard Form 990 is used by larger nonprofits, Form 990-EZ is a shorter version for smaller organizations, and Form 990-PF applies to private foundations.} I first extract these fields directly from structured XML nodes, allowing for variation in tag names across filing formats. For address information, I accommodate both U.S. and foreign address structures, selecting the appropriate format based on the presence of relevant XML nodes. When structured fields are missing or incomplete, I implement a text-based fallback procedure. I read the raw XML content as text and use pattern matching to recover organization name, EIN, location, and filing year from header-style blocks in the filing. This approach substantially improves coverage in cases where XML structure is inconsistent. I validate each firm observation by checking for missing identifiers, names, location, or year, and flag observations requiring review. I ultimately restrict the sample to U.S.-based organizations. Specifically, I drop filings that explicitly report a foreign country in the address fields and remove all associated person-level observations for those firm-year combinations.

    \subsubsection{Person-Level Data}

    I construct a person-level dataset capturing leadership team members listed in each filing. For each individual, I extract name, title, average hours worked per week, role indicators (e.g., officer, director, trustee, key employee), and compensation (from the organization and related entities). I identify individuals by searching across multiple XML node structures, including Form 990 Part VII, Form 990-EZ officer sections, and Form 990-PF equivalents. To ensure robustness across filing formats, I also implement a generic search for XML nodes that simultaneously contain name and title fields. If no structured nodes are detected, I implement a text-based fallback procedure. This procedure scans the raw filing text to identify uppercase name strings and matches them to nearby lines containing recognized leadership titles (e.g., CEO, CFO, director, trustee). This step recovers leadership information from filings with non-standard formatting. I remove duplicate entries within each filing and retain only observations with non-missing names.


    \subsection{Data Cleaning}

    I begin by dropping invalid placeholder names and removing observations that reflect compensation descriptions rather than individuals. I convert all text to lowercase and remove duplicate records. I then clean and standardize the name field and construct role-related indicators. I create a cleaned name variable by removing punctuation, digits, parenthetical information, and references to schedules or compensation notes. Using manually curated lists of credentials and keywords (e.g., MD, RN, MBA), I generate indicator variables capturing whether an individual is a physician, nurse, or healthcare administrator. I also strip these credentials and other extraneous descriptors from the name field to isolate the underlying person name. To address inconsistencies in how names are recorded within the same organization, I implement three complementary matching procedures at the EIN level. First, I use approximate string matching (based on edit distance) to reconcile minor spelling differences, keeping the most frequently occurring version of a name. Second, I construct a graph of names that share at least two tokens (e.g., first and last name) and collapse them to a canonical version (the shortest string) within each connected component. Third, I address formatting inconsistencies such as missing spaces (e.g., “johnsmith” versus “john smith”) by identifying prefix matches and again consolidating to a canonical form. Together, these steps standardize names within organizations and reduce duplication arising from formatting variation or minor errors.

    Next, I extract first names for use in gender prediction. I begin with a Social Security–based list of common first names (augmented with additional manually identified names) and tokenize each cleaned name into its component words. I use a series of decision rules, applied at the EIN-year level, to infer which token most likely represents the first name. When the majority of individuals in an organization follow a “first-name-first” convention, I assign the first token accordingly; when instead names frequently appear in “initial-lastname” or alternative formats, I assign the second token under additional conditions. I then fill in missing assignments by leveraging consistency within the same name and organization across observations, and apply additional row-level rules based on whether exactly one token matches a known first name or whether specific common names (e.g., “mary” or “john”) appear in the string. Manual checks confirm that these procedures yield consistent first-name assignments across years for the same individual.

    I assign organizational roles using a structured combination of title text, name text, and the role indicator variables available in the XML filings. I first assign core executive roles using keyword matching, including CEO, CFO, COO, CMO, CIO, CNO, vice president (VP), and other executive positions. These assignments rely on common variants and frequent misspellings (e.g., “chief executive,” “c.e.o.,” “chief medical officer”) and are implemented sequentially; once an individual is assigned to one of these roles, that assignment is not overwritten by later rules.
    
    I then identify additional leadership roles that may not correspond to senior executives, including administrators, departmental leaders, and individuals with “director of” or similar titles. Individuals with titles indicating divisional or functional leadership (e.g., “division chief,” “medical director,” “director of…”) are classified as non-executive leaders when they have not already been assigned an executive role. I identify board members using both the Form 990 indicator variables and title-based rules. Individuals flagged exclusively as trustees or directors (and not as officers, key employees, or highly compensated employees) are classified as board members regardless of their title. I also classify individuals as board members when they hold officer status but have titles consistent with governance roles (e.g., chair, secretary, treasurer, trustee, or director). Additional title-based rules are used to capture board members when indicator variables appear incomplete or inconsistently reported. To separate executives from highly compensated clinical staff, I classify individuals with clinical occupations (e.g., physicians, nurses, and other medical specialists) as employees when they have not already been assigned to an executive or leadership role. Remaining individuals with “director” titles and substantial reported hours are classified as non-executive leaders, while all unassigned individuals—who are overwhelmingly specialized clinical staff—are classified as employees.

    A key challenge arises with individuals whose titles include “president”, as they can either represent the president of the hospital or of the board. I initially classify such individuals as potential executive presidents when they receive positive compensation, and then resolve their role using additional information. First, individuals whose titles include terms such as “foundation,” “auxiliary,” or “department” are reclassified as non-executive leaders. Second, when a CEO is already identified within a hospital-year, remaining presidents are treated as board members. Third, when no CEO is recorded, I promote a potential executive president to CEO if that individual appears as CEO in adjacent years or if they are the only presidential candidate in a hospital-year that otherwise contains executive officers. When multiple presidential candidates exist in the absence of a CEO, I rank them based on indicators of executive responsibility including key employee status, hospital-related title keywords, and reported hours, and assign the highest-ranked individual as CEO, classifying the remaining candidates as board members. Finally, when no other executive roles are observed in a hospital-year, presidential titles are treated as board roles, reflecting cases where executive leadership is not reported in the filing.
        
    After these steps, I make additional adjustments to ensure consistency across hospital-years. I correct likely classification errors by leveraging within-person consistency over time when a role appears to be a one-year deviation from persistent assignments, and I resolve cases with multiple CEOs per hospital-year by prioritizing non-system-level executives when both system and hospital CEOs are present. I then restrict the sample to individuals with meaningful time commitment to the organization, dropping individuals whose reported hours suggest they primarily work for related entities rather than the focal hospital. The final dataset consists of executive team members—including CEO, CFO, COO, CMO, CIO, CNO, and other senior executives whose roles are consistently defined across hospital-years.

    I classify executives as having a clinical background using information from both reported titles and occupation descriptions. Specifically, I identify individuals as clinical when their titles indicate medical training or clinical practice, including physicians, surgeons, nurses, anesthesiologists, and other licensed healthcare professionals. These classifications rely on keyword matching applied to the reported title field, capturing a broad set of clinical roles commonly observed in Form 990 filings (e.g., "md", “physician,”, “registered nurse,” “nurse practitioner,” “anesthesiologist”). If an individual is ever classified as clinical, I code them as always clinical. Additionally, I assume that all Chief Medical Officers and Chief Nursing Officers are clinical. 

    I make several limitations to the person-level data. First, I remove individuals that are former personnel based on the form's ``former" indicator, and by checking if the word former is in the name or title. Second, I classify executives as full-time, part-time, or effectively external (``no-time") based on whether at least 75\% of their reported hours are worked at the specified organization versus related organizations, and then, within each firm-year, I drop ``no-time” executives when there are more than a minimal number of actively working leaders. Finally, if the hospital distinguishes between system versus hospital executives, I drop system-level individuals.

    \subsection{Aggregation to the Hospital Level}

    I construct hospital-year level measures of executive team composition, size, and turnover by aggregating individual-level classifications from Form 990 filings. For each hospital and year, I compute team size as the number of individuals classified as executives. I then measure composition using the share of executives who are female and the share with clinical backgrounds. I also compute compensation measures, including total CEO compensation and average compensation per executive, using both organization-level and related-organization compensation fields. I treat observations with zero reported compensation as missing, as these values likely reflect reporting error or incomplete reporting rather than true compensation. I also construct measures of turnover and team dynamics at the hospital-year level, including entries and exits, and a measure of overall team stability. 

    \textit{Here I will write about merging the hospital level data with other data sets}

    \textit{Here I will talk about any limitations I make}

    




    \section{Summary Statistics}

    In Table \ref{tab:exec_table}, I summarize key variables relating to the hospital executive teams. On average, executive teams are small, with a mean size of just over \textcolor{red}{N} members, though there is substantial variation across hospitals. Executive teams include both clinical and non-clinical leaders: clinicians account for roughly \textcolor{red}{N} percent of executives, while women represent about \textcolor{red}{N} percent on average, indicating meaningful but incomplete representation along both dimensions. Compensation levels vary widely, with mean CEO compensation of \textcolor{red}{N} and substantial dispersion. In any given year, \textcolor{red}{N} percent of hospitals experience some entry into the executive team, while CEO and CMO transitions occur in roughly \textcolor{red}{N} percent of hospital-years. I define team stability as the Jaccard similarity between consecutive executive teams, equal to the number of individuals present in both years divided by the total number of distinct individuals present in either year. Teams exhibit a moderate degree of stability, with approximately \textcolor{red}{N} percent of executives remaining from one year to the next. Consistent with this pattern, entry and exit counts are typically low in any given year, but the combination of frequent turnover events and relatively stable team cores suggests that changes in executive composition occur through gradual and ongoing adjustments rather than large, discrete shifts. The number of observations varies across variables; in particular, CMO-related measures are based on fewer observations because not all hospitals have a Chief Medical Officer at any given point in time. 


    \import{Objects}{summary_table.tex}


    Additionally, I present average values over time in Figure \ref{fig:hospital_exec_figure}. Compensation rises steadily over the period, with CEO compensation increasing from roughly \textcolor{red}{N} to \textcolor{red}{N}, while compensation per executive follows a similar but more muted upward trajectory. At the same time, representation within executive teams shifts gradually: the share of female executives and those with clinical backgrounds both trend upward, though these changes are modest. The turnover panel shows that executive teams experience frequent but relatively stable churn, with approximately \textcolor{red}{N} percent of hospitals experiencing some entry in a given year and CEO and CMO transitions occurring in about \textcolor{red}{N} percent of hospital-years. Despite this ongoing turnover, overall team stability remains high—around \textcolor{red}{N} percent, indicating that most executives persist from one year to the next. Finally, team size increases slightly over the sample period, rising from \textcolor{red}{N} executives to around \textcolor{red}{N} on average. Taken together, these patterns suggest that changes in executive composition occur through gradual adjustments within largely stable teams, rather than through large, discrete shifts in leadership structure. 


    \begin{figure}[ht!]
        \centering
        \caption{Executive Team Characteristics Over Time}
        \includegraphics[width=\linewidth]{Objects/hosp_avgs_time.pdf}
        \label{fig:hospital_exec_figure}
    \end{figure}

    key takeaways I can use from this graph: panel 1 -> "growth in clinical presence coming more from the team than the CEO position panel 2 -> "clinical leadership is not just a byproduct of generic turnover" and panel 3 -> "most CEO turnover does not result in a clinical CEO" 

    



    %\import{Objects}{composition_dynamics_table.tex}









    \section{Correlation with Hospital Outcomes}

    \section{Conclusion}
    
    



    

    

	
	
	


\end{document}